\section{Geometric, algebraic, and encoding conventions}

\subsection{Finite triangulations and local identifications}

Let $T$ be a finite generalized triangulation with $t$ tetrahedra. Each
tetrahedron has four numbered vertices. A face is either unpaired or is
paired to another face by a permutation of the four local vertices.
The reciprocal pairing and its inverse are both recorded. Local corners,
edges, and faces are distinguished from global equivalence classes under
these identifications. In particular, several corners of a tetrahedron
can represent the same global vertex.

The implemented geometric validator requires a connected orientable
compact three-manifold with one torus boundary component. It checks
reciprocal pairings, orientation consistency, edge identifications,
vertex links, and the boundary's connectedness and Euler characteristic.
An interior vertex link is a sphere; a boundary vertex link is a disc.
These are the hypotheses under which the following normal-coordinate and
Pachner operations are used. Establishing that an arbitrary supplied
triangulation is the exterior of the user's knot is a separate provenance
obligation, supported by the repository's diagram-exterior routines.

\subsection{Normal coordinates}

Write a standard normal vector as
\[
 x=(a_{iv},q_{ij})\in\R_{\ge0}^{7t},
 \qquad 0\le i<t,\quad 0\le v<4,\quad 0\le j<3.
\]
Here $a_{iv}$ counts triangles cutting off vertex $v$ in tetrahedron $i$,
and $q_{ij}$ counts quadrilaterals of type $j$. Across each paired face,
the three normal-arc counts agree. These conditions form a sparse integer
linear system $Ax=0$. Each arc equation is a difference of two triangle
coordinates and two quadrilateral coordinates, with coincidences combined.
The additional quadrilateral condition requires at most one positive
$q_{ij}$ in each tetrahedron. It is this nonconvex condition that necessitates
sector choices or a different global method.

On the linear matching space the Euler characteristic is a rational
linear functional $c^Tx$. One may calculate it by assigning each global
edge-intersection count to a representative local edge and counting one
copy of each glued normal arc. For integral admissible vectors the result
is the integer Euler characteristic of the corresponding surface.
The linear formula remains algebraically defined on relaxed vectors
with incompatible quadrilaterals; those vectors do not necessarily
represent embedded surfaces.

The full integral vector can be very large in value but small in encoding.
Let $B$ bound the bit length of its coordinates or local height differences.
Algorithms here manipulate $O(t)$ binary integers, not one object per
normal disc. Any theoretical polynomial LP bound refers to the full
rational input bit length. The shipped exact simplex is an implementation
choice with explicit resource limits, and has no polynomial pivot bound.

\subsection{Anchors and vertex links}

For a global vertex $v$, let $C_v$ be its set of local corner records.
The normal vertex-link vector $L_v$ has one triangle at every corner in
$C_v$ and zero quadrilaterals. The maximal removable multiplicity is
\[
 m_v(x)=\min_{c\in C_v} a_c,
 \qquad
 \peel(x)=x-\sum_v m_v(x)L_v.
\]
The peeled vector is nonnegative, obeys the same matching equations,
and is admissible when $x$ is. Choosing one triangle coordinate at each
global vertex and setting it to zero gives an \emph{anchor face}. Every
peeled vector lies on some anchor face. A one-vertex triangulation has
at most $4t$ raw corner anchors, with additional contractions possible
as in the incoming dual report.

Peeling is distinct from division by the common divisor of quadrilateral
coordinates. The latter can be useful for component-count queries, but
it changes all coordinates and is not part of the local-update theorem.
For $\lambda_v=\chi(L_v)\in\{1,2\}$ and $d_v=|C_v|$,
\begin{equation}\label{eq:peel-identities}
 \chi(\peel(x))=\chi(x)-\sum_v\lambda_vm_v(x),
 \qquad
 P(\peel(x))=P(x)-\sum_vd_vm_v(x),
\end{equation}
where $P$ is the number of normal discs before gluing. These identities
also explain why vertices shared by disjoint local replacements can
couple their peeled scores.

\subsection{Integral cocycles and coherent fibres}

An integral cocycle can be represented in each tetrahedron by four
integer heights $h_{i0},\ldots,h_{i3}$. The signed height differences
agree on globally identified edges. Adding the same constant to all
four heights of one tetrahedron does not change those differences.
Sorting its four heights as $h_a\le h_b\le h_c\le h_d$ gives the
coherent normal vector with $h_b-h_a$ triangles at $a$,
$h_d-h_c$ triangles at $d$, and $h_c-h_b$ quadrilaterals separating
$\{a,b\}$ from $\{c,d\}$. Equal heights cause zero counts and require
no special perturbation.

The construction is the normal level set of the local affine height
functions, with unit-separated levels taken modulo integral shifts.
It is cooriented and its boundary and normal matching data agree under
face gluings. A transport operation preserves the integral cohomology
class but can split, join, or change components of this fibre. The
incoming transport package contains a verified example in which one
disc becomes a disc and two spheres \cite{priorTransport}. Thus component
analysis remains necessary after any geometric search policy.

\subsection{Cost model and certificate scope}

The comparison model counts arithmetic on integers separately from its
bit cost. An operation on a bounded number of $B$-bit heights takes
polynomial bit time in $B$; multiplication by a corner count introduces
$O(\log t)$ additional bits. We write $\M(b)$ for a chosen integer
multiplication cost and avoid treating a large integer as a machine word.
Combinatorial indices use $O(\log(t+2))$ bits.

Three kinds of output are deliberately distinct: an algebraic certificate
about a supplied matching system; a geometric certificate relating two
triangulations and transported cocycles; and a knot verdict tied to the
original diagram. The first two can accelerate the third but do not
replace its end-to-end proof requirements.
